% GENERATED FILE. Edit canonical lessons, then run npm run generate:books.
% canonical-source-hash: fnv1a64:75827f1a439e6f05
% canonical-lessons: JA-C394-sakamichi, JA-C394-kewashii, JA-C394-haba, JA-C394-kudaru, JA-C394-koeru

\chapter{Up the Hill and Down}
\label{ch:ja-up-the-hill-and-down}
\hlchaptermodality{394}

\begin{chapteropening}
\textbf{By the end of this chapter:} \emph{I can describe steep roads and mountain crossings.}

Complete the last lesson of chapter 394: I can describe steep roads and mountain crossings.
\end{chapteropening}

\section[\texorpdfstring{\ja{さかみち} (sakamichi)}{さかみち (sakamichi)}]{\ja{さかみち} (sakamichi) — a hill road; a sloping street}

\label{lesson:JA-C394-sakamichi}



\begin{quote}
Before the new one: say the Japanese for an arrow (on a sign), then the Japanese for to look back.
\end{quote}

\subsection*{You'll want to know: \ja{さかみち}}
\textbf{\ja{さかみち}} — \emph{sakamichi} — \textquotedblleft{}a hill road; a sloping street\textquotedblright{}.

\begin{quote}
\textbf{\ja{さかみちを} \ja{ゆっくり} \ja{のぼりました。}} — \emph{sakamichi o yukkuri noborimashita} — \textquotedblleft{}I slowly climbed the hill road.\textquotedblright{}
\end{quote}

\begin{cousinweb}
\textbf{\ja{さかみち}} (\emph{sakamichi}) is a native Japanese word (wago), inherited rather than borrowed from Chinese, built from \emph{saka}, \textquotedblleft{}slope\textquotedblright{} and \emph{michi}, \textquotedblleft{}road\textquotedblright{}; its kanji were attached later as kun'yomi (JMdict).
\end{cousinweb}

\begin{grammarlens}[title={What you've built}]
One more word for getting around.
\end{grammarlens}

\subsection*{Your turn}
\begin{itemize}
\raggedright
  \item \emph{Say it:} \emph{sakamichi}
  \item \emph{Say it:} \emph{sakamichi}, once more
  \item \emph{Say it:} \emph{sakamichi o yukkuri noborimashita} — I slowly climbed the hill road.
  \item \emph{Recall:} say the Japanese for an arrow (on a sign), then the Japanese for to look back, then say \emph{sakamichi} again
\end{itemize}

\begin{tcolorbox}[breakable,colback=teal!4,colframe=teal!35!black,title={Before you move on}]
What does \ja{さかみち} mean? (\textquotedblleft{}A hill road; a sloping street\textquotedblright{}.) Say it once more.
\end{tcolorbox}
\section[\texorpdfstring{\ja{けわしい} (kewashii)}{けわしい (kewashii)}]{\ja{けわしい} (kewashii) — steep; rugged}

\label{lesson:JA-C394-kewashii}



\begin{quote}
Before the new one: say the Japanese for to look back, then the Japanese for a hill road.
\end{quote}

\subsection*{You'll want to know: \ja{けわしい}}
\textbf{\ja{けわしい}} — \emph{kewashii} — \textquotedblleft{}steep; rugged\textquotedblright{}.

\begin{quote}
\textbf{\ja{やまの} \ja{みちは} \ja{けわしいです。}} — \emph{yama no michi wa kewashii desu} — \textquotedblleft{}The mountain path is steep.\textquotedblright{}
\end{quote}

An i-adjective.

\begin{cousinweb}
\textbf{\ja{けわしい}} (\emph{kewashii}) is a native Japanese word (wago), inherited rather than borrowed from Chinese; its kanji was attached later as a kun'yomi (JMdict).
\end{cousinweb}

\begin{grammarlens}[title={What you've built}]
One more word for getting around.
\end{grammarlens}

\subsection*{Your turn}
\begin{itemize}
\raggedright
  \item \emph{Say it:} \emph{kewashii}
  \item \emph{Say it:} \emph{kewashii}, once more
  \item \emph{Say it:} \emph{yama no michi wa kewashii desu} — The mountain path is steep.
  \item \emph{Recall:} say the Japanese for to look back, then the Japanese for a hill road, then say \emph{kewashii} again
\end{itemize}

\begin{tcolorbox}[breakable,colback=teal!4,colframe=teal!35!black,title={Before you move on}]
What does \ja{けわしい} mean? (\textquotedblleft{}Steep; rugged\textquotedblright{}.) Say it once more.
\end{tcolorbox}
\section[\texorpdfstring{\ja{はば} (haba)}{はば (haba)}]{\ja{はば} (haba) — width}

\label{lesson:JA-C394-haba}



\begin{quote}
Before the new one: say the Japanese for a hill road, then the Japanese for steep.
\end{quote}

\subsection*{You'll want to know: \ja{はば}}
\textbf{\ja{はば}} — \emph{haba} — \textquotedblleft{}width\textquotedblright{}.

\begin{quote}
\textbf{\ja{この} \ja{みちは} \ja{はばが} \ja{せまいです。}} — \emph{kono michi wa haba ga semai desu} — \textquotedblleft{}This road is narrow.\textquotedblright{}
\end{quote}

\begin{cousinweb}
\textbf{\ja{はば}} (\emph{haba}) is a native Japanese word (wago), inherited rather than borrowed from Chinese; its kanji was attached later as a kun'yomi (JMdict).
\end{cousinweb}

\begin{grammarlens}[title={What you've built}]
One more word for getting around.
\end{grammarlens}

\subsection*{Your turn}
\begin{itemize}
\raggedright
  \item \emph{Say it:} \emph{haba}
  \item \emph{Say it:} \emph{haba}, once more
  \item \emph{Say it:} \emph{kono michi wa haba ga semai desu} — This road is narrow.
  \item \emph{Recall:} say the Japanese for a hill road, then the Japanese for steep, then say \emph{haba} again
\end{itemize}

\begin{tcolorbox}[breakable,colback=teal!4,colframe=teal!35!black,title={Before you move on}]
What does \ja{はば} mean? (\textquotedblleft{}Width\textquotedblright{}.) Say it once more.
\end{tcolorbox}
\section[\texorpdfstring{\ja{くだる} (kudaru)}{くだる (kudaru)}]{\ja{くだる} (kudaru) — to descend; to go down (a slope or river)}

\label{lesson:JA-C394-kudaru}



\begin{quote}
Before the new one: say the Japanese for steep, then the Japanese for width.
\end{quote}

\subsection*{You'll want to know: \ja{くだる}}
\textbf{\ja{くだる}} — \emph{kudaru} — \textquotedblleft{}to descend; to go down (a slope or river)\textquotedblright{}.

\begin{quote}
\textbf{\ja{さかみちを} \ja{くだって、かわまで} \ja{いきました。}} — \emph{sakamichi o kudatte, kawa made ikimashita} — \textquotedblleft{}I went down the hill road as far as the river.\textquotedblright{}
\end{quote}

Polite \textbf{\ja{くだります}}, \emph{kudarimasu}; a godan verb. Past (polite): \textbf{\ja{ふねが} \ja{かわを} \ja{くだりました。}}, \emph{fune ga kawa o kudarimashita}.

\begin{cousinweb}
\textbf{\ja{くだる}} (\emph{kudaru}) is a native Japanese word (wago), inherited rather than borrowed from Chinese; its kanji was attached later as a kun'yomi (JMdict).
\end{cousinweb}

\begin{grammarlens}[title={What you've built}]
One more word for getting around.
\end{grammarlens}

\subsection*{Your turn}
\begin{itemize}
\raggedright
  \item \emph{Say it:} \emph{kudaru}
  \item \emph{Say it:} \emph{kudaru}, once more
  \item \emph{Say it:} \emph{sakamichi o kudatte, kawa made ikimashita} — I went down the hill road as far as the river.
  \item \emph{Recall:} say the Japanese for steep, then the Japanese for width, then say \emph{kudaru} again
\end{itemize}

\begin{tcolorbox}[breakable,colback=teal!4,colframe=teal!35!black,title={Before you move on}]
What does \ja{くだる} mean? (\textquotedblleft{}To descend; to go down (a slope or river)\textquotedblright{}.) Say it once more.
\end{tcolorbox}
\section[\texorpdfstring{\ja{こえる} (koeru)}{こえる (koeru)}]{\ja{こえる} (koeru) — to cross over (a mountain or border)}

\label{lesson:JA-C394-koeru}



\begin{quote}
Before the new one: say the Japanese for width, then the Japanese for to descend.
\end{quote}

\subsection*{You'll want to know: \ja{こえる}}
\textbf{\ja{こえる}} — \emph{koeru} — \textquotedblleft{}to cross over (a mountain or border)\textquotedblright{}.

\begin{quote}
\textbf{\ja{やまを} \ja{こえて、となりの} \ja{まちに} \ja{いきました。}} — \emph{yama o koete, tonari no machi ni ikimashita} — \textquotedblleft{}I crossed the mountain and went to the next town.\textquotedblright{}
\end{quote}

Polite \textbf{\ja{こえます}}, \emph{koemasu}; an ichidan verb. Past (polite): \textbf{\ja{ちちが} \ja{やまを} \ja{こえました。}}, \emph{chichi ga yama o koemashita}.

\begin{cousinweb}
\textbf{\ja{こえる}} (\emph{koeru}) is a native Japanese word (wago), inherited rather than borrowed from Chinese; its kanji was attached later as a kun'yomi (JMdict).
\end{cousinweb}

\begin{grammarlens}[title={What you've built}]
One more word for getting around.
\end{grammarlens}

\subsection*{Your turn}
\begin{itemize}
\raggedright
  \item \emph{Say it:} \emph{koeru}
  \item \emph{Say it:} \emph{koeru}, once more
  \item \emph{Say it:} \emph{yama o koete, tonari no machi ni ikimashita} — I crossed the mountain and went to the next town.
  \item \emph{Recall:} say the Japanese for width, then the Japanese for to descend, then say \emph{koeru} again
\end{itemize}

\begin{tcolorbox}[breakable,colback=teal!4,colframe=teal!35!black,title={Before you move on}]
What does \ja{こえる} mean? (\textquotedblleft{}To cross over (a mountain or border)\textquotedblright{}.) Say it once more.
\end{tcolorbox}
